\documentclass[10pt,a4paper]{article}
\usepackage[utf8]{inputenc}
\usepackage{amsmath}
\usepackage{amsfonts}
\usepackage{amssymb}
\usepackage{graphicx}
\usepackage{float}
\usepackage{booktabs}
\usepackage{enumerate}
\title{GSF-3100 \\
06 Marché obligataire gouvernemental \\
Série d'exercices 5}

\begin{document}
\maketitle

\begin{itemize}
\item Dans ce document vous avez des questions à choix multiples allant de a) à d).  
\item Il y a seulement un choix possible par question.
\item Ce questionnaire contient des exercices en lien avec le marché obligataire gouvernemental et plus pécisément sur le marché primaire et l'adjudiction d'obligations.
\item Contient 6 questions. 
\end{itemize}
 
\newpage

\section*{Mise en contexte}

Le gouvernement fédéral annonce une nouvelle émission d’obligations standards d’une valeur totale de 20 millions \$.  En réponse, il reçoit les offres de deux distributeurs détaillées ci-bas.  

\begin{table}[H]
\centering
\resizebox{\textwidth}{!}{%
\begin{tabular}{@{}llll@{}}
\toprule
Types d'offres & Offres non concurentielles & Offres concurrentielles #1 & Offres concurrentielles #2 \\ \midrule
Distributeur A & 4 millions \$               & 1 millions \$ à y=2\%        & 1 millions \$ à y=3\%        \\
Distributeur B & 11 millions \$              & 2 millions \$ à y=4\%        & 1 millions \$ à y=5\%        \\ \bottomrule
\end{tabular}%
}
\end{table}

\section*{Exercice 1}
Déterminer en millions de dollars,  le montant total des offres non concurrentielles acceptées.

\begin{enumerate}[a)]
\item 4 millions \$
\item 11 millions \$
\item 15 millions \$
\item 20 millions \$
\end{enumerate}

\section*{Exercice 2}
Déterminer en millions de dollars,  le montant total des offres concurrentielles acceptées.

\begin{enumerate}[a)]
\item 4 millions \$
\item 5 millions \$
\item 11 millions \$
\item 15 millions \$
\end{enumerate}

\section*{Exercice 3}
Déterminer le taux de rendement des offres non concurrentielles si l'émission est effectuée par le gouvernement canadien.

\begin{enumerate}[a)]
\item 3.00\%
\item 3.50\%
\item 3.60\%
\item 4.00\%
\end{enumerate}

\section*{Exercice 4}
Déterminer le taux de coupon des obligations émises si l'émission est effectuée par le gouvernement canadien.

\begin{enumerate}[a)]
\item 3.00\%
\item 3.50\%
\item 3.60\%
\item 4.00\%
\end{enumerate}

\section*{Exercice 5}
Déterminer le taux de rendement des offres non concurrentielles si l'émission est effectuée par par le gouvernement américain.

\begin{enumerate}[a)]
\item 3.50\%
\item 4.00\%
\item 4.50\%
\item 5.00\%
\end{enumerate}

\section*{Exercice 6}
Déterminer le taux de coupon des obligations émises si l'émission est effectuée par par le gouvernement américain.

\begin{enumerate}[a)]
\item 4.50\%
\item 4.75\%
\item 5.00\%
\item 5.25\%
\end{enumerate}


\end{document}
